% GENERATED FILE. Edit canonical lessons, then run npm run generate:books.
% canonical-source-hash: fnv1a64:614d2eb07a85e9c9
% canonical-lessons: JA-C173-tabemono, JA-C173-nomimono, JA-C173-obentou, JA-C173-nokoru, JA-C173-ajiwau

\chapter{Food, Drinks and Leftovers}
\label{ch:ja-food-drinks-and-leftovers}
\hlchaptermodality{173}

\begin{chapteropening}
\textbf{By the end of this chapter:} \emph{I can talk about food and drink, and what is left over.}

Complete the last lesson of chapter 173: I can talk about food and drink, and what is left over.
\end{chapteropening}

\section[\texorpdfstring{\ja{たべもの} (tabemono)}{たべもの (tabemono)}]{\ja{たべもの} (tabemono) — food}

\label{lesson:JA-C173-tabemono}



\begin{quote}
Before the new one: say the Japanese for dinner, then the Japanese for to grill.
\end{quote}

\subsection*{You'll want to know: \ja{たべもの}}
\textbf{\ja{たべもの}} — \emph{tabemono} — \textquotedblleft{}food\textquotedblright{}.

\begin{quote}
\textbf{\ja{にほんの} \ja{たべものは} \ja{おいしいです。}} — \emph{nihon no tabemono wa oishii desu} — \textquotedblleft{}Japanese food is tasty.\textquotedblright{}
\end{quote}

\begin{cousinweb}
\textbf{\ja{たべもの}} (\emph{tabemono}) is native Japanese, built from \emph{tabe}, the stem of \emph{taberu}, \textquotedblleft{}to eat\textquotedblright{}, and \emph{mono}, \textquotedblleft{}thing\textquotedblright{} (JMdict, KANJIDIC2).
\end{cousinweb}

\begin{grammarlens}[title={What you've built}]
One more word for everyday life at A2.
\end{grammarlens}

\subsection*{Your turn}
\begin{itemize}
\raggedright
  \item \emph{Say it:} \emph{tabemono}
  \item \emph{Say it:} \emph{tabemono}, once more
  \item \emph{Say it:} \emph{nihon no tabemono wa oishii desu} — Japanese food is tasty.
  \item \emph{Recall:} say the Japanese for dinner, then the Japanese for to grill, then say \emph{tabemono} again
\end{itemize}

\begin{tcolorbox}[breakable,colback=teal!4,colframe=teal!35!black,title={Before you move on}]
What does \ja{たべもの} mean? (\textquotedblleft{}Food\textquotedblright{}.) Say it once more.
\end{tcolorbox}
\section[\texorpdfstring{\ja{のみもの} (nomimono)}{のみもの (nomimono)}]{\ja{のみもの} (nomimono) — beverage; something to drink}

\label{lesson:JA-C173-nomimono}



\begin{quote}
Before the new one: say the Japanese for to grill, then the Japanese for food.
\end{quote}

\subsection*{You'll want to know: \ja{のみもの}}
\textbf{\ja{のみもの}} — \emph{nomimono} — \textquotedblleft{}beverage; something to drink\textquotedblright{}.

\begin{quote}
\textbf{\ja{のみものは} \ja{おちゃを} \ja{たのみます。}} — \emph{nomimono wa ocha o tanomimasu} — \textquotedblleft{}For a drink, I order tea.\textquotedblright{}
\end{quote}

\begin{cousinweb}
\textbf{\ja{のみもの}} (\emph{nomimono}) is native Japanese, built from \emph{nomi}, the stem of \emph{nomu}, \textquotedblleft{}to drink\textquotedblright{}, and \emph{mono}, \textquotedblleft{}thing\textquotedblright{} (JMdict, KANJIDIC2).
\end{cousinweb}

\begin{grammarlens}[title={What you've built}]
One more word for everyday life at A2.
\end{grammarlens}

\subsection*{Your turn}
\begin{itemize}
\raggedright
  \item \emph{Say it:} \emph{nomimono}
  \item \emph{Say it:} \emph{nomimono}, once more
  \item \emph{Say it:} \emph{nomimono wa ocha o tanomimasu} — For a drink, I order tea.
  \item \emph{Recall:} say the Japanese for to grill, then the Japanese for food, then say \emph{nomimono} again
\end{itemize}

\begin{tcolorbox}[breakable,colback=teal!4,colframe=teal!35!black,title={Before you move on}]
What does \ja{のみもの} mean? (\textquotedblleft{}Beverage; something to drink\textquotedblright{}.) Say it once more.
\end{tcolorbox}
\section[\texorpdfstring{\ja{おべんとう} (obentō)}{おべんとう (obentō)}]{\ja{おべんとう} (obentō) — boxed lunch; bento}

\label{lesson:JA-C173-obentou}



\begin{quote}
Before the new one: say the Japanese for food, then the Japanese for beverage.
\end{quote}

\subsection*{You'll want to know: \ja{おべんとう}}
\textbf{\ja{おべんとう}} — \emph{obentō} — \textquotedblleft{}boxed lunch; bento\textquotedblright{}.

\begin{quote}
\textbf{\ja{こうえんで} \ja{おべんとうを} \ja{たべます。}} — \emph{kōen de obentō o tabemasu} — \textquotedblleft{}I eat a boxed lunch in the park.\textquotedblright{}
\end{quote}

\begin{cousinweb}
\textbf{\ja{おべんとう}} (\emph{obentō}) is built from polite \emph{o-} and Sino-Japanese \emph{bentō}, two characters read with their on'yomi, \emph{ben} and \emph{tō}, readings borrowed from Chinese; the older history of the compound is not settled (JMdict, KANJIDIC2).
\end{cousinweb}

\begin{grammarlens}[title={What you've built}]
One more word for everyday life at A2.
\end{grammarlens}

\subsection*{Your turn}
\begin{itemize}
\raggedright
  \item \emph{Say it:} \emph{obentō}
  \item \emph{Say it:} \emph{obentō}, once more
  \item \emph{Say it:} \emph{kōen de obentō o tabemasu} — I eat a boxed lunch in the park.
  \item \emph{Recall:} say the Japanese for food, then the Japanese for beverage, then say \emph{obentō} again
\end{itemize}

\begin{tcolorbox}[breakable,colback=teal!4,colframe=teal!35!black,title={Before you move on}]
What does \ja{おべんとう} mean? (\textquotedblleft{}Boxed lunch; bento\textquotedblright{}.) Say it once more.
\end{tcolorbox}
\section[\texorpdfstring{\ja{のこる} (nokoru)}{のこる (nokoru)}]{\ja{のこる} (nokoru) — to be left over; to remain}

\label{lesson:JA-C173-nokoru}



\begin{quote}
Before the new one: say the Japanese for beverage, then the Japanese for boxed lunch.
\end{quote}

\subsection*{You'll want to know: \ja{のこる}}
\textbf{\ja{のこる}} — \emph{nokoru} — \textquotedblleft{}to be left over; to remain\textquotedblright{}.

\begin{quote}
\textbf{\ja{ごはんが} \ja{すこし} \ja{のこります。}} — \emph{gohan ga sukoshi nokorimasu} — \textquotedblleft{}A little rice is left over.\textquotedblright{}
\end{quote}

Polite \textbf{\ja{のこります}}, \emph{nokorimasu}; a godan verb. Its transitive partner \textbf{\ja{のこす}} is taught already.

\begin{cousinweb}
\textbf{\ja{のこる}} (\emph{nokoru}) is a native Japanese word (wago), inherited rather than borrowed from Chinese; its kanji was attached later as a kun'yomi, and it shares its stem with \emph{nokosu}, \textquotedblleft{}to leave behind\textquotedblright{} (JMdict, KANJIDIC2).
\end{cousinweb}

\begin{grammarlens}[title={What you've built}]
One more word for everyday life at A2.
\end{grammarlens}

\subsection*{Your turn}
\begin{itemize}
\raggedright
  \item \emph{Say it:} \emph{nokoru}
  \item \emph{Say it:} \emph{nokoru}, once more
  \item \emph{Say it:} \emph{gohan ga sukoshi nokorimasu} — A little rice is left over.
  \item \emph{Recall:} say the Japanese for beverage, then the Japanese for boxed lunch, then say \emph{nokoru} again
\end{itemize}

\begin{tcolorbox}[breakable,colback=teal!4,colframe=teal!35!black,title={Before you move on}]
What does \ja{のこる} mean? (\textquotedblleft{}To be left over; to remain\textquotedblright{}.) Say it once more.
\end{tcolorbox}
\section[\texorpdfstring{\ja{あじわう} (ajiwau)}{あじわう (ajiwau)}]{\ja{あじわう} (ajiwau) — to savour; to taste, to relish}

\label{lesson:JA-C173-ajiwau}



\begin{quote}
Before the new one: say the Japanese for boxed lunch, then the Japanese for to be left over.
\end{quote}

\subsection*{You'll want to know: \ja{あじわう}}
\textbf{\ja{あじわう}} — \emph{ajiwau} — \textquotedblleft{}to savour; to taste, to relish\textquotedblright{}.

\begin{quote}
\textbf{\ja{ゆっくり} \ja{おちゃを} \ja{あじわいます。}} — \emph{yukkuri ocha o ajiwaimasu} — \textquotedblleft{}I slowly savour my tea.\textquotedblright{}
\end{quote}

Polite \textbf{\ja{あじわいます}}, \emph{ajiwaimasu}; a godan verb.

\begin{cousinweb}
\textbf{\ja{あじわう}} (\emph{ajiwau}) is a native Japanese verb, inherited, built on \emph{aji}, \textquotedblleft{}taste\textquotedblright{}; its kanji was attached later as a kun'yomi (JMdict, KANJIDIC2).
\end{cousinweb}

\begin{grammarlens}[title={What you've built}]
One more word for everyday life at A2.
\end{grammarlens}

\subsection*{Your turn}
\begin{itemize}
\raggedright
  \item \emph{Say it:} \emph{ajiwau}
  \item \emph{Say it:} \emph{ajiwau}, once more
  \item \emph{Say it:} \emph{yukkuri ocha o ajiwaimasu} — I slowly savour my tea.
  \item \emph{Recall:} say the Japanese for boxed lunch, then the Japanese for to be left over, then say \emph{ajiwau} again
\end{itemize}

\begin{tcolorbox}[breakable,colback=teal!4,colframe=teal!35!black,title={Before you move on}]
What does \ja{あじわう} mean? (\textquotedblleft{}To savour; to taste, to relish\textquotedblright{}.) Say it once more.
\end{tcolorbox}
